\documentclass[11pt]{article}
\usepackage[margin=1in]{geometry}
\usepackage{parskip}
\usepackage{hyperref}
\hypersetup{colorlinks=true, urlcolor=blue}
\usepackage{url}
\usepackage{xurl}
\usepackage{xcolor}

\begin{document}

\begin{center}
{\LARGE\bfseries Title Page}
\end{center}

\vspace{1em}
\noindent\textbf{Manuscript title:}\\
Counterfactual Attribution of Crypto Systemic Stress on a Ledger-Specified Causal Graph: Methods and a Cautionary Benchmark

\vspace{1em}
\noindent\textbf{Authors (in submission-system order):}\\
1. Quang-Vinh Dang --- British University Vietnam, Hung Yen, Vietnam

\vspace{1em}
\noindent\textbf{Corresponding author:}\\
Quang-Vinh Dang, British University Vietnam, Hung Yen, Vietnam.
E-mail: \href{mailto:vinh.dq4@buv.edu.vn}{vinh.dq4@buv.edu.vn}

\bigskip
\hrulefill

\vspace{1.5em}
\begin{center}
{\LARGE\bfseries Author's Information}
\end{center}

\vspace{1em}
\noindent\textbf{Quang-Vinh Dang.}\quad
Quang-Vinh Dang is a Lecturer at British University Vietnam and Data/AI
Director at SmartOSC, where he leads initiatives in multi-agent AI systems,
large language models, and graph neural networks across banking, retail, and
finance. He holds a Ph.D. in Computer Science from Universit\'e de Lorraine
(LORIA), France (2018), an M.Sc. in Financial Engineering from WorldQuant
University (2021), and completed a Post-Master programme in Software
Technology at Eindhoven University of Technology. His research spans machine
learning in finance, natural language processing, intrusion detection, and
federated learning, and he has authored more than 60 peer-reviewed papers
indexed in Scopus/SCI-E, many in Q1 journals. Prior to his current roles, he
held data-leadership positions at Trusting Social, Be Group, GoTymeX, and
DatVietVAC, and was a postdoctoral researcher at Inria, France.

\bigskip
\hrulefill

\vspace{1.5em}
\begin{center}
{\LARGE\bfseries Cover Letter}
\end{center}

\vspace{1em}
\noindent Quang-Vinh Dang\\
British University Vietnam\\
Hung Yen, Vietnam\\
\href{mailto:vinh.dq4@buv.edu.vn}{vinh.dq4@buv.edu.vn}

\vspace{1em}
\noindent The Editor-in-Chief\\
\emph{Digital Finance}

\vspace{1em}
\noindent Dear Editor-in-Chief,

\vspace{0.5em}
I am submitting the manuscript \textbf{``Counterfactual Attribution of Crypto
Systemic Stress on a Ledger-Specified Causal Graph: Methods and a Cautionary
Benchmark''} for consideration in \emph{Digital Finance}. Measures of crypto
systemic risk usually infer contagion networks from price co-movement, using
CoVaR-type or expected-shortfall statistics. The resulting graphs are estimated
rather than observed, and the attributions built on them cannot be checked
against the transfers that moved funds between entities. Public blockchains
record those transfers, but the raw ledger becomes an entity-level graph only
after addresses are resolved to institutions. The October 2025 liquidation
cascade made this gap concrete.

The manuscript treats the entity-resolved on-chain ledger as the candidate parent
set of a linear structural causal model of entity stress. Its methodological
contribution is a framework that is analytic throughout: an interventional tail
measure (On-Chain Causal CoVaR) and a counterfactual loss attribution computed by
full abduction, with flow-free network-multiplier bounds, abduction invariance of
stress deviations, a share bound, a walk representation of the linear Shapley
attribution with an exact $O(2^p)$ subset algorithm, and a certified ranking
radius. The paper also shows that the choice of abduction changes the answer. On
the same data, zero-shock abduction inflates the Binance-to-Bybit share, so
attributions built this way need full abduction to be interpretable.

The paper is also a cautionary benchmark, and I present it as one. Out of
sample, the network equation explains little hourly stress: the pooled
out-of-sample $R^2$ is 0.005, against 0.161 for an own-lag model.
No channel's share exceeds 15.2\%. H1, that the observed graph predicts
better than a price-estimated graph, is not supported, and edge direction is not
distinguishable from reversal. Placebo evidence is mixed. I argue that this is
the value of the paper for supervisors and risk teams: a ledger-specified graph
is auditable, and auditability is not the same as explanatory power.

The paper fits the scope of \emph{Digital Finance} on three counts. It addresses
systemic risk in crypto-asset markets and how that risk is measured. It addresses
supervisory transparency, since each attribution traces back to public ledger
records and an explicit model. It also provides reproducible inputs: the code, the
query definitions, and the derived result files are openly available, and the
core data come from free public sources.

This manuscript is original, has not been published previously, and is not under
consideration at any other journal. It is single-authored, and the author declares
no competing interests. The code and derived results are publicly available at
\url{https://github.com/vinhqdang/VEIN-On-Chain-Causal-Risk-Attribution-for-Cryptocurrency-Systemic-Risk}.
Thank you for considering this manuscript.

\vspace{1.5em}
\noindent Sincerely,

\vspace{0.5em}
\noindent Quang-Vinh Dang\\
British University Vietnam

\bigskip
\hrulefill

\vspace{1.5em}
\begin{center}
{\LARGE\bfseries Declarations}
\end{center}

\vspace{1em}
\noindent\textbf{Ethical Approval and Consent to participate.} Not
applicable: this study uses only public, aggregated on-chain and market data
and involves no human subjects, human data, or animals.

\vspace{0.8em}
\noindent\textbf{Consent for publication.} Not applicable.

\vspace{0.8em}
\noindent\textbf{Availability of supporting data.} All code, query definitions and
derived result files underlying every number, table, and figure are publicly
available at
\url{https://github.com/vinhqdang/VEIN-On-Chain-Causal-Risk-Attribution-for-Cryptocurrency-Systemic-Risk}.
Raw transfers are regenerated from the public BigQuery dataset
\texttt{bigquery-public-data.crypto\_ethereum} (token transfers), which can be
queried in the free BigQuery sandbox. Address labels come from the eth-labels and
etherscan-labels repositories, both MIT-licensed, and a label table built from
these two sources alone is included. Prices are from CoinGecko and protocol TVL
from DefiLlama. No restricted or proprietary data were used.

\vspace{0.8em}
\noindent\textbf{Competing interests.} The author declares no competing
interests.

\vspace{0.8em}
\noindent\textbf{Funding.} No external funding supported this work.

\vspace{0.8em}
\noindent\textbf{Authors' contributions.} Q.-V.D. conceived the study,
developed the VEIN framework and its implementation, performed the data
collection and empirical analysis, and wrote the manuscript.

\vspace{0.8em}
\noindent\textbf{Acknowledgements.} Not applicable.

\vspace{0.8em}
\noindent\textbf{Authors' information (Optional).} See ``Author's
Information'' above.

\end{document}
